% !TeX root = main_long.tex
\section{RQ4: comparison with existing semantic models}
\subsection{Method}
\label{sec:m-rq4}

% AUTHOR REVIEW R20: identify who assigned the RQ4 judgments; the available record does not name this person.
Published artifacts for GFVO, HERO-Genomics, GVO, and VCF2RDF were inspected against nine retrieval criteria (Table~\ref{tab:models})~\cite{baran2015gfvo,cazzaro2025hero,kawashima2023gvo,penha2017vcf2rdf}.
Verdicts are \textbf{E}, explicitly represented; \textbf{C}, representable with additional conventions; and \textbf{S}, outside stated scope.
The criteria concern VCF interpretation context, and S denotes a difference of purpose rather than a defect.
F7 covers only per-allele lists: \texttt{Number=G} values index possible genotypes, and although VCF Core records each item's genotype index, decoding it into alleles relies on the specification's ordering and is untested (requirement R37).
Artifact digests identify the inspected versions. 
VCF2RDF's terms were inspected, while converter behaviour was taken from its paper rather than execution.
GA4GH VRS is not scored because it addresses computable variation identity rather than source-file representation~\cite{wagner2021vrs}.
This assessment measures documented support, not the capabilities of extensions or applications built around these models.
The verdicts reflect one inspection, without an inter-rater assessment.

\subsection{Results}
\label{sec:models}

Table~\ref{tab:models} summarizes the inspected artifacts.
HERO-Genomics and GVO retain INFO as one literal through \term{hero:vcfInfo} and \term{gvo:info}, requiring consumers to parse the text and supply interpretation rules~\cite{cazzaro2025hero,kawashima2023gvo}.
GFVO directly models field meanings such as \term{AlleleCount} and \term{MappingQuality}, but leaves the declaration mechanism outside its scope~\cite{baran2015gfvo}.
VCF2RDF publishes declaration terms, including \term{Number}, \term{Type}, and \term{Description}, but no allele-index term, so allele association requires additional conventions~\cite{penha2017vcf2rdf}.
GVO's predominantly S verdicts reflect its focus on classifying biological variation rather than representing source-file relationships.

\begin{table}[htbp]
\caption{Support in inspected artifacts: E, explicitly represented; C, requiring additional conventions; S, outside stated scope. These are inspection judgments about VCF interpretation context.}
\label{tab:models}
\centering\small
\begin{tabularx}{\linewidth}{@{}Xccccc@{}}
\toprule
Can a user\ldots & VCF Core & GFVO & HERO & GVO & VCF2RDF \\
\midrule
F1\quad trace a record to its file & E & E & E & S & C \\
F2\quad read declaration ID, Number, and Type & E & S & S & S & E \\
F3\quad read the declared VCF version & E & S & S & S & C \\
F4\quad address ALT allele $n$ by index & E & C & C & S & C \\
F5\quad separate ploidy, allele calls, and phasing & E & C & C & S & C \\
F6\quad get a (sample, field) value unsplit & E & C & C & S & C \\
F7\quad bind a \texttt{Number=A/R} item to its allele & E & C & C & S & C \\
F8\quad distinguish explicit \texttt{.} from unstated & E & S & C & C & C \\
F9\quad identify the sample encoding & E & S & S & S & S \\
\bottomrule
\end{tabularx}
\end{table}

VCF Core's distinguishing role is the explicit representation of source interpretation context: declarations, version rules, allele indices, and sample scope.
These findings concern the published artifacts and do not show that other models cannot support the same information through extensions.
Separately, VCF Core provides \textbf{123 curated alignments}~\footnote{\url{https://github.com/ecrum19/vcf-core-vocabulary/tree/main/mappings}}, recorded in SSSOM with SKOS predicates: 57 exact, 40 close, 16 related, and ten narrow matches~\cite{matentzoglu2022sssom,skos2009}.
These graded term mappings support connections to other models without establishing identity between instances.
